\documentclass[11pt,a4paper]{article}
\usepackage[margin=24mm,headheight=14pt]{geometry}
\usepackage[T1]{fontenc}
\usepackage{lmodern,microtype}
\usepackage{amsmath,amssymb,amsthm,mathtools}
\usepackage{booktabs,tabularx,array,enumitem}
\usepackage[dvipsnames]{xcolor}
\usepackage{xurl}
\usepackage[colorlinks=true,linkcolor=MidnightBlue,citecolor=MidnightBlue,urlcolor=MidnightBlue]{hyperref}
\usepackage{fancyhdr}
\pagestyle{fancy}
\fancyhf{}
\fancyhead[L]{\small BundleFuse: score sufficiency}
\fancyhead[R]{\small Theory and evidence}
\fancyfoot[C]{\thepage}
\renewcommand{\headrulewidth}{0.3pt}
\setlength{\parindent}{0pt}
\setlength{\parskip}{6pt plus 1pt}
\setlength{\emergencystretch}{2em}
\clubpenalty=10000
\widowpenalty=10000
\setlist[itemize]{leftmargin=1.5em,itemsep=2pt,topsep=3pt}
\setlist[enumerate]{leftmargin=1.7em,itemsep=2pt,topsep=3pt}
\newtheorem{proposition}{Proposition}
\newtheorem{definition}{Definition}
\newcommand{\E}{\mathbb E}
\newcommand{\R}{\mathbb R}
\newcommand{\ind}{\mathbf 1}
\newcommand{\norm}[1]{\left\lVert #1\right\rVert}
\newcommand{\off}{\operatorname{off}}
\newcommand{\diag}{\operatorname{diag}}
\newcommand{\Normal}{\mathcal N}
\newcommand{\rownorm}{\operatorname{rownorm}}
\newcolumntype{Y}{>{\raggedright\arraybackslash}X}
\hypersetup{pdftitle={When Are BundleFuse's Scores Sufficient?},pdfauthor={BundleFuse research analysis},pdfsubject={Estimator structure, sufficiency, representational limits, and a stopping rule}}

\begin{document}
\thispagestyle{plain}
{\Large\bfseries When Are BundleFuse's Scores Sufficient?\par}
{\large A theory of what to retain, what can be proved, and where to stop\par}
{\small Standalone analysis \hfill 7 November 2026\par}
\vspace{3pt}\hrule\vspace{7pt}

\textbf{The question.} BundleFuse can construct several scores from historical
bundles and item attributes. Why retain these scores, and why not add another
one? A satisfactory answer must explain the estimation problem, identify what
an extra score could contribute, and specify a defensible stopping condition.
Good benchmark results alone do not determine a necessary number of scores.

\textbf{The conclusion.} BundleFuse is a complete estimator whose final rule
has an exact optimization interpretation. Its organizing idea is to start from
relationships shared across items and permit departures supported by a
candidate's history. The present score list is not universally sufficient:
an explicit four-item construction defeats it. A rigorous, achievable target is
a compact estimator that is within a stated ranking tolerance of a declared
family of alternatives under a resource budget.

This document derives the relevant statements and distinguishes their status:
an identity of the implemented estimator, a theorem under stated assumptions,
a verified counterexample, or a proposed statistical certificate. The
foundational sufficiency and shrinkage arguments are established mathematics;
their application here does not establish that these principles are new.

\section{The prediction problem and the source of the difficulty}
\label{sec:problem}

Let $\mathcal I$ be a catalogue of $n$ items. A query supplies a nonempty
observed set $S\subset\mathcal I$. The eligible candidates are
$\mathcal I_S=\mathcal I\setminus S$. We must rank them to recover a nonempty
hidden set $G\subseteq\mathcal I_S$. The training collection contains $m$
bundles, represented by the binary matrix
\begin{equation}
 B\in\{0,1\}^{m\times n},\qquad
 B_{bi}=1\ \Longleftrightarrow\ \text{training bundle $b$ contains item $i$}.
\end{equation}
Define the item frequency $f_i=\sum_b B_{bi}$ and the bundle size
$d_b=\sum_i B_{bi}$. These definitions concern binary, unit-weight bundles,
as in the fashion analysis below. Other data profiles may use weighted inputs.

Two shoes can share a category and have similar descriptions. Such shared
attributes support a starting prediction for both. If only one repeatedly
appears with the particular trousers and shirt in $S$, recorded item identities
can distinguish them. For a newly introduced shoe, however, no historical
match could mean either an unsuitable combination or simply an unobserved
combination. The zero count alone cannot decide which.

This motivates one estimation problem: \emph{use relationships that can be
shared across items, then allow a candidate's specific observations to adjust
the prediction}. A candidate with little recorded history should have a
restricted ability to override that starting point. Candidate frequency is an
available measure of exposure; it is not a calibrated measure of correctness,
and frequent appearance does not imply frequent appearance with this query.

The different score constructions are possible ways of estimating these
relationships. Their names do not imply independent information, and a count
of score columns is not a count of independent observations.

\newpage
\section{One estimator: a shared prediction and a controlled departure}
\label{sec:estimator}

For a fixed candidate $c$ and query $S$, let $Y_c$ denote the selected full
score and $P_c$ the selected restricted fallback. Let $s>0$ be the
validation-frozen fallback scale, and define $F_c=sP_c$. The selected
fallback strength is $\kappa\geq0$. For $\kappa>0$, BundleFuse computes
\begin{equation}
 \boxed{\quad
 \theta_c=F_c+\rho_c(Y_c-F_c),\qquad
 \rho_c=\frac{f_c}{f_c+\kappa}.
 \quad}
 \label{eq:core}
\end{equation}
All these scores depend on $S$; the query argument is suppressed for
readability. The difference $Y_c-F_c$ is the proposed departure from the
fallback. The coefficient $\rho_c$ determines how much survives.
For $f_c=0$, the result is $F_c$. Holding $Y_c$ and $F_c$ fixed, the result
approaches $Y_c$ as $f_c$ increases.

The fallback's information must be described accurately. Category relations
pool different item IDs carrying the same label. Description scores compare
attributes. The bridge also uses historical neighbors, and the fallback can
include popularity. Thus $F_c$ is neither an independent estimate nor an
exclusively descriptive one. Since $s$ need not equal one, the departure
$Y_c-F_c$ need not contain only historical terms.

\subsection{An exact optimization interpretation}

\begin{proposition}[The implemented interpolation]
For fixed $Y_c,F_c,f_c$ and $\kappa>0$, equation~\eqref{eq:core} is the
unique minimizer over real $t$ of
\begin{equation}
 J_c(t)=\frac{f_c}{2}(t-Y_c)^2+\frac{\kappa}{2}(t-F_c)^2.
 \label{eq:objective}
\end{equation}
\end{proposition}
\begin{proof}
The derivative is
$J'_c(t)=(f_c+\kappa)t-f_cY_c-\kappa F_c$.
Its zero gives equation~\eqref{eq:core}. The second derivative is
$f_c+\kappa>0$, so the minimizer is unique.
\end{proof}

The two terms charge for departing from the full score and from the fallback,
respectively. The first charge grows with historical exposure. This
interpretation needs no probability model for the scores. It does not derive
the raw score menu or establish that its weights are optimal.

When $\kappa=0$, the implementation returns $Y_c$. If also $f_c=0$, the
objective in~\eqref{eq:objective} is constant and does not uniquely specify
that answer. The explicit convention is therefore part of the estimator.

There is also an exact stability property. If $Y_c$ changes by $e_c$ while
$F_c$ stays fixed, $\theta_c$ changes by $\rho_c e_c$. For candidates
$c$ and $d$, the change in their score difference is at most
\begin{equation}
 \rho_c|e_c|+\rho_d|e_d|.
\end{equation}
This limits sensitivity to changes confined to the full score, especially
at low frequency. It does not guarantee better rankings: a suppressed
departure could have been correct. If a new source changes the fallback too,
that change must also be included.

\subsection{What the dense model actually selects}

Let $r_h(c\mid S)$ be a raw score with source index $h$. In the examined
full-ranking implementation it is standardized over all $n$ catalogue items
before observed items are masked:
\begin{align}
 \mu_h(S)&=\frac1n\sum_i r_h(i\mid S), &
 \sigma_h(S)^2&=\frac1{n-1}\sum_i[r_h(i\mid S)-\mu_h(S)]^2,\notag\\
 z_h(c\mid S)&=\frac{r_h(c\mid S)-\mu_h(S)}{\sigma_h(S)+10^{-9}}.
 \label{eq:standardize}
\end{align}
The candidate comparisons use different functions of the same available
data. Standardization balances their numerical scales; it does not establish
independence or calibrate them as probabilities.

For the frozen POG\_DENSE research recipe, define $Q_c=P_c$. The selected
construction is
\begin{align}
 Q_c&=0.5z_{\rm cat}(c)+0.25z_{\rm bridge}(c)+0.25z_{\rm pop}(c),\notag\\
 Y_c&=Q_c+0.5z_{\rm simp}(c)+0.25z_{\rm csvd}(c),\notag\\
 F_c&=sQ_c,\qquad s\simeq1.46055159,\qquad \kappa=0.5,\notag\\
 \theta_c&=[s+(1-s)\rho_c]Q_c
       +\rho_c[0.5z_{\rm simp}(c)+0.25z_{\rm csvd}(c)].
 \label{eq:dense}
\end{align}

\begin{center}
\begin{tabularx}{\linewidth}{@{}lY@{}}
\toprule
Selected score & Relationship being measured\\
\midrule
cat & Training relations between item categories, with a penalty for repeating categories that rarely repeat within a bundle.\\
bridge & Similarity between category-centered descriptors after incorporating descriptions of historical neighbors.\\
pop & The candidate's log training frequency, $\log(1+f_c)$.\\
simp & Candidate support from individual query items and from query-item pairs appearing together with it.\\
csvd & Similarity between compressed profiles of item co-occurrence associations.\\
\bottomrule
\end{tabularx}
\end{center}

Appendix~\ref{app:scores} defines these five raw scores mathematically.
The dense recipe does not select standalone cooc, ppmi, soft-group,
propagated-score or Rao columns. PPMI and propagation still appear inside
selected representations; their absence as separate columns does not remove
those internal operations. The frozen POG recipe uses four columns; the
Spotify recipe uses seven weighted columns plus its explicit Rao term.
None of these counts has been proved necessary or minimal.

\section{What information is sufficient for a ranking?}
\label{sec:ranking}

A sufficiency claim must identify the decision and the population. Let $X$
contain all permitted prediction-time information: $S$, the training data,
item attributes, candidate eligibility and any allowed user context. The
hidden set $G$ is random under the specified population of queries. Let
$K\leq|\mathcal I_S|$ be the cutoff and let $\pi(j)$ be the candidate at
position $j$.

For Recall@$K$, define the candidate contribution $V_c$ and position weight
$\omega_j$ by
\begin{equation}
 V_c=\frac{\ind\{c\in G\}}{|G|},\qquad \omega_j=1\quad(1\leq j\leq K).
 \label{eq:recalltarget}
\end{equation}
For binary-relevance NDCG@$K$, instead define
\begin{equation}
 \omega_j=\frac1{\log_2(j+1)},\qquad
 D_K(G)=\sum_{j=1}^{\min(K,|G|)}\omega_j,\qquad
 V_c=\frac{\ind\{c\in G\}}{D_K(G)}.
 \label{eq:ndcgtarget}
\end{equation}
Here $D_K(G)$ is the largest attainable discounted gain for that hidden set.
In either case, realized utility equals $\sum_{j=1}^K\omega_jV_{\pi(j)}$.
Define the expected contribution
\begin{equation}
 \eta_c(X)=\E[V_c\mid X].
\end{equation}
The conditional expected utility of a ranking is therefore
\begin{equation}
 U(\pi\mid X)=\sum_{j=1}^K\omega_j\eta_{\pi(j)}(X).
 \label{eq:expectedutility}
\end{equation}

Sorting $\eta$ maximizes this expression. To see this, place candidates $c$
and $d$ at positions $j<k$, where $\omega_j\geq\omega_k$. If
$\eta_c<\eta_d$, exchanging them improves utility by
$(\omega_j-\omega_k)(\eta_d-\eta_c)\geq0$. Treat positions below $K$ as
having weight zero. Repeated exchanges produce a sorted ranking.
For Recall only the selected set matters; for NDCG its order matters too.

Thus a representation need only preserve an optimal ordering, rather than
the entire distribution of $G$. When $|G|$ varies, plain membership
probabilities need not have the same ordering as these expected normalized
contributions. The NDCG distinction is established in the consistency
analysis of Ravikumar, Tewari and Yang~\cite{ravikumar2011}.

BundleFuse's numerical scores are not calibrated estimates of $\eta$.
Consequently, errors between its raw numerical $\theta$ and $\eta$ cannot
be interpreted as calibrated prediction errors without further justification.
Only their rankings are directly compared in the existing evaluation.

\section{The value of one more score}
\label{sec:additional}

Let $Z$ retain the full vectors of selected candidate scores, frequencies and
eligibility information. Let $R$ be a proposed additional score vector. These
are functions of allowed inputs, not hidden labels. For the contribution
vector $V$ defined for the chosen metric, write
\begin{align}
 m_0(Z)&=\E[V\mid Z], & m_1(Z,R)&=\E[V\mid Z,R],\notag\\
 \Delta(R\mid Z)&=\E\big[\norm{m_1-m_0}_2^2\big].
 \label{eq:delta}
\end{align}
The norm sums over eligible candidates. Expectations refer to the declared
query population. $\Delta$ measures how much the best prediction of the
relevant items changes after observing $R$. It is not a correlation between
raw scores and not the standalone performance of $R$.

\begin{proposition}[Exact decrease in optimal prediction error]
\begin{equation}
 \E\big[\norm{V-m_0}_2^2\big]
 -\E\big[\norm{V-m_1}_2^2\big]=\Delta(R\mid Z).
 \label{eq:projection}
\end{equation}
\end{proposition}
\begin{proof}
Expand $V-m_0=(V-m_1)+(m_1-m_0)$. Conditional on $(Z,R)$,
$V-m_1$ has mean zero, while $m_1-m_0$ is fixed. The expected cross term
is therefore zero. Expanding the squared norm gives the identity.
\end{proof}

If $\Delta=0$, revealing $R$ cannot improve optimal expected ranking utility
for this metric: the two conditional mean vectors are equal almost surely.
A positive $\Delta$ need not improve utility, because changed values can
retain the same top-$K$ ordering. This result uses the classical projection
property of conditional expectation~\cite{banerjee2005}. Its relation to
statistical sufficiency is that relevant prediction information is preserved;
full conditional independence is a stronger condition~\cite{fukumizu2004}.

\begin{proposition}[Bound on additional ranking utility]
Let $\boldsymbol\omega=(\omega_1,\ldots,\omega_K)$. The improvement in
optimal expected utility from revealing $R$ is at most
\begin{equation}
 \min\left\{1,\ 2\norm{\boldsymbol\omega}_2
                    \sqrt{\Delta(R\mid Z)}\right\}.
 \label{eq:gainbound}
\end{equation}
\end{proposition}
\begin{proof}
Represent a ranking by a vector $a$ assigning weight $\omega_j$ to its
$j$th candidate and zero below $K$. Every ranking vector has norm
$\norm{\boldsymbol\omega}_2$. Let $a_0$ maximize $a^\top m_0$ and
$a_1$ maximize $a^\top m_1$. Then
\begin{align}
 (a_1-a_0)^\top m_1
 &= (a_1-a_0)^\top m_0+(a_1-a_0)^\top(m_1-m_0)\notag\\
 &\leq 2\norm{\boldsymbol\omega}_2\norm{m_1-m_0}_2.
\end{align}
The first term is nonpositive by optimality of $a_0$; the second is bounded
by the Cauchy--Schwarz and triangle inequalities. Take expectations and
use $\E[\norm{W}_2]\leq\sqrt{\E[\norm{W}_2^2]}$. Both evaluation
metrics lie in $[0,1]$, which supplies the cap of one.
\end{proof}

For a specified class of admissible extensions, a population-level stopping
condition is therefore
\begin{equation}
 \sup_R\Delta(R\mid Z)\leq
 \frac{\varepsilon^2}{4\norm{\boldsymbol\omega}_2^2}.
 \label{eq:populationstop}
\end{equation}
Under this condition, no admissible extension improves the best $Z$-based
ranking by more than $\varepsilon$. The class and population must be
specified; the inequality is not known to hold for the present BundleFuse
scores. It cannot simply be assumed as a proof of their sufficiency.

\subsection{Existing information can also be used better}

The current combiner might not extract all the information in $Z$. Let
$h(Z)$ be any numerical prediction on the same scale as $V$, and define
$E_0=\E[\norm{h-m_0}_2^2]$. A second orthogonality identity gives
\begin{equation}
 \E[\norm{h-m_1}_2^2]=E_0+\Delta(R\mid Z).
 \label{eq:extraction}
\end{equation}
Indeed, $\E[m_1-m_0\mid Z]=0$, making its expected inner product with
$h-m_0$ zero. The same ranking argument bounds the gain over the ranking
of $h$ by $2\norm{\boldsymbol\omega}_2\sqrt{E_0+\Delta}$, capped at one.
A calibrated $h$ with BundleFuse's ranking has not been established, so this
is a conditional interpretation, not a measured certificate for $\theta$.

A new model can thus help by extracting existing information better, even
when it adds no information. PPMI and SVD being derived from the same
matrix does not make their score columns interchangeable for a restricted
weighted combiner.

These results also apply when $R$ jointly contains several new scores.
Testing additions individually cannot rule out joint benefit. Two independent
fair bits each reveal nothing about their exclusive-or, while together they
determine it. Likewise, a nonsignificant addition does not prove
$\Delta=0$; the experiment may simply be too imprecise.

\section{Why the present score menu is not universally complete}
\label{sec:counterexample}

An explicit construction shows what the scores can miss. Create four item
groups $A,B,C,D$, each containing $q$ values indexed $0,\ldots,q-1$. For
a fixed shift $h\in\{0,\ldots,q-1\}$, include every training bundle
\begin{equation}
 \{A_a,B_b,C_c,D_{(a+b+c+h)\bmod q}\},
 \qquad a,b,c=0,\ldots,q-1.
 \label{eq:corpus}
\end{equation}
There are $q^3$ bundles, each of size four. Use the group as the category
label and give all items within each group the same descriptor.

Fixing $r$ values from distinct groups, for $r=1,2,3$, leaves exactly
$q^{3-r}$ possible bundles, independently of $h$. A subset containing two
different values from one group occurs zero times. Hence complete
single-item, pair and triple counts are identical for all shifts. All bundle
size discounts are identical as well.

For $S=\{A_0,B_0,C_0\}$, however, the recorded completion is $D_h$.
The four-item count identifies it immediately. With $q=2$, the two possible
shifts give the particularly small example
\begin{center}
\begin{tabular}{@{}ccc@{}}
\toprule
Observed triple & Completion when $h=0$ & Completion when $h=1$\\
\midrule
$\{A_0,B_0,C_0\}$ & $D_0$ & $D_1$\\
\bottomrule
\end{tabular}
\end{center}

The failure of the examined fashion scores is specific:
\begin{itemize}
\item Direct pair, co-occurrence and PPMI scores use identical frequency
and pair statistics. PPMI-SVD receives the same matrix, with identical
association-profile rows for all $D$ candidates.
\item The incidence correction uses candidate--seed--seed triple counts,
which are also identical. Its derivation appears in
Appendix~\ref{app:history}.
\item Propagation uses the same pair operator. Group-identical descriptors
cannot distinguish $D$ candidates; their category-centered residuals are zero.
\item Category, soft-group and popularity scores cannot distinguish these
candidates either. Rao uses the same candidate--seed affinities and the
same relations among the query items.
\end{itemize}

Consequently, standardization and frequency interpolation receive equal
inputs for every $D$ value and produce equal scores in ideal arithmetic.
An ID-based tie-break does not reconstruct the four-item relation. With
$q=21$, more $D$ candidates are tied than there are top-20 positions.
Varying $h$ makes any one of them the correct answer while preserving the
score inputs, so no fixed top-20 rule from these scores succeeds for every
corpus. The ordering limitation for NDCG already appears with two values.

The construction was evaluated with the production fashion score operators
and the frozen dense fusion for $q=2$ and $q=21$. It verified equal counts
through order three and score equality to floating-point tolerance. The
largest score differences between the two checked shifts were
$2.39\times10^{-7}$ and $1.20\times10^{-7}$, respectively.
These are synthetic mathematical checks, not benchmark performance results.

The conclusion is limited but decisive: a theorem that this menu captures
every possible completion relationship is false. It does not follow that
such four-item relationships are frequent or useful in the intended data.
That becomes an empirical question with a precise missing relationship to
test, rather than an invitation to add arbitrary scores indefinitely.

\section{A stopping rule with a finite statistical guarantee}
\label{sec:stopping}

Condition~\eqref{eq:populationstop} is conceptually useful but generally
difficult to establish. A more practical standard compares a finite,
predeclared family of fitted systems under the ranking metric itself.

Let $\mathcal H$ contain $M$ systems whose fitted parameters and procedures
are frozen before a fresh audit sample is examined. Include the compact
candidate, relevant source removals, close replacements and justified
additions. Let $C(h)$ be a defined cost, such as inference time under a fixed
measurement protocol. Memory and fitting cost can be separate hard
constraints. Let $\varepsilon\geq0$ be the acceptable absolute utility loss.

Suppose the audit contains $N$ independent units from the target population.
Each system's utility on each unit lies in $[0,1]$. A unit can be a complete
original bundle with a predefined average across its query masks. Repeated
masks of the same bundle must not be counted as independent observations.
User or temporal dependence can require a different unit or probability
bound. Let $U(h)$ be population utility and $\widehat U(h)$ its audit mean.
Condition on the independently fitted systems when applying the following
argument.

For a failure probability $\delta\in(0,1)$, define
\begin{equation}
 b=\sqrt{\frac{\log(2M/\delta)}{2N}}.
 \label{eq:hoeffding}
\end{equation}
Hoeffding's inequality~\cite{hoeffding1963} and a union bound imply
\begin{equation}
 \Pr\left\{\text{for every }h\in\mathcal H,
       \ |\widehat U(h)-U(h)|\leq b\right\}\geq1-\delta.
\end{equation}

\begin{proposition}[A compact-system guarantee]
Choose a least-cost member $h_{\rm small}$ satisfying
\begin{equation}
 \widehat U(h_{\rm small})\geq
       \max_{h\in\mathcal H}\widehat U(h)-\varepsilon.
\end{equation}
With probability at least $1-\delta$,
\begin{equation}
 U(h_{\rm small})\geq
       \max_{h\in\mathcal H}U(h)-\varepsilon-2b.
 \label{eq:certificate}
\end{equation}
\end{proposition}
\begin{proof}
On the simultaneous event, the selected system's population utility is at
least its audit utility minus $b$. Its audit utility is at least the best
audit utility minus $\varepsilon$. The best audit utility is at least the
population-best member's utility minus $b$. Combine the inequalities.
\end{proof}

Simultaneity permits selecting among these frozen systems after seeing the
audit means. It does not permit fitting new systems on the audit labels and
pretending the original guarantee still applies. A particular baseline is
certified within tolerance $\varepsilon$ of the family if
\begin{equation}
 \max_{h\in\mathcal H}
    [\widehat U(h)-\widehat U(\text{baseline})]+2b\leq\varepsilon.
\end{equation}

The elementary bound is often loose. With $M=100$, $N=4{,}000$ and
$\delta=0.05$, $2b\simeq0.0644$ absolute utility. That exceeds several
observed differences in this project. A narrow certificate would need a
justified tighter paired bound or more independent data. A nonsignificant
comparison is not a substitute for such a bound.

This defines adequacy relative to a population, metric, tolerance, resource
constraint and comparison family. It does not cover an algorithm invented
after the audit. It does provide a reason to stop: retain a compact system
once justified alternatives cannot offer enough gain to warrant their cost.

\section{What can be claimed about BundleFuse now?}
\label{sec:claims}

The completed experiments support particular choices rather than an
indispensable universal menu. The controlled incidence bank improves dense
POG NDCG over its pairwise counterpart. On POG, history and shared relations
work together and the fallback helps. On Spotify, a different historical
overlap estimator improved the controlled result. Clothing and sparse
Spotify discarded the explicit historical column in the follow-up.
These findings favor a compact, dataset-selected implementation.

They do not establish the certificate in Section~\ref{sec:stopping}.
The studies are exploratory supplied-split comparisons, conditional on
earlier source inventories, and their test outcomes have informed the
research. They did not search every subset or audit all relevant joint
additions on fresh data.

A further caution concerns the normal observation model derived in
Appendix~\ref{app:normal}. It provides a real sufficient-aggregate theorem
under its assumptions, but those assumptions are not established for the
current standardized scores. The saved full recipes fit the mean squared
$(Y-F)$ against inverse frequency and obtain these raw slopes:
\begin{center}
\begin{tabular}{@{}lrr@{}}
\toprule
Recipe & Raw inverse-frequency slope & Selected $\kappa$\\
\midrule
POG & $-0.273238$ & $0.5$\\
POG\_DENSE & $-2.528107$ & $0.5$\\
Spotify & $-5.335897$ & $5.0$\\
\bottomrule
\end{tabular}
\end{center}
These estimates do not substantiate a positive inverse-frequency noise
relationship. They are not themselves formal rejection tests. The procedure
truncates inadmissible variance estimates and considers validation-selected
alternatives. Thus equation~\eqref{eq:objective} is an unconditional account
of the current interpolation; a calibrated posterior is not established.

\subsection{A publication statement with a definite object of study}

\begin{quote}
BundleFuse estimates how an item completes an observed set from recorded
item combinations and relationships shared across items. Its historical
coefficients are determined by bundle counts. Its final rule limits
departures from a shared prediction according to candidate history. The
scorer requires no gradient-trained scoring network; its scalar choices are
selected on validation data.
\end{quote}

The exact historical estimator, the chosen sharing operations, their
computation, and their controlled empirical effects are the objects to
defend. The supplied item features can have upstream training costs; SVD,
clustering and scalar selection are also fitted operations. ``No
gradient-trained scoring network'' is more precise than ``nothing is
trained.''

The reported comparisons with published CLHE/RaMen values support empirical
effectiveness under the project's documented comparison scope. They do not
prove score sufficiency, identify the cause of the advantage, or establish
priority for the full construction. RaMen's task is bundle
construction~\cite{ramen2025}. Related bundle recommendation already uses
frequency-dependent mixtures, for example CoHeat~\cite{coheat2024},
although its task and learned representations differ. A mathematical
property and a claim of originality require separate evidence.

The response to ``why not add another model?'' should be precise:
\begin{quote}
BundleFuse specifies a complete estimator with a limited fitting and
computational budget. Retained components must justify their incremental
ranking value. An extension is relevant when it addresses a missing
relationship or improves the accuracy--cost tradeoff. The contribution does
not depend on claiming that no future extension can ever help.
\end{quote}

The present analysis establishes exact behavior and representational limits.
It supplies a route to a compact-system guarantee. It does not certify that
the current number of scores is minimal, or turn established sufficiency
theorems into new algorithmic contributions.

\appendix
\section{Definitions of the five dense score columns}
\label{app:scores}

This appendix makes the dense example self-contained. It defines the
mathematical operations used in equation~\eqref{eq:dense}. Apart from the
standardization stabilizer in~\eqref{eq:standardize}, tiny numerical
stabilizers in logs and row norms are omitted for readability. A zero vector
is mapped to zero by row normalization. No proof here depends on those
floating-point conventions.

\subsection{Individual-seed and joint-seed historical support}
\label{app:history}

Define the seed weight and bundle weight
\begin{equation}
 a_i=(f_i+5)^{-1/2},\qquad v_b=(d_b+1)^{-1/2}.
\end{equation}
For a query, let $o_b(S)=\sum_{s\in S}B_{bs}a_s$ be its weighted overlap
with training bundle $b$, and $k_b(S)=\sum_{s\in S}B_{bs}$ its number of
overlapping items. The direct and additional scores are
\begin{align}
 p_c(S)&=\sum_b B_{bc}v_b o_b(S),\notag\\
 \delta_c(S)&=\beta\sum_b B_{bc}v_b o_b(S)[k_b(S)-1]_+,
       \qquad \beta=1.5,\notag\\
 g_c(S)&=p_c(S)+\delta_c(S),\qquad r_{\rm simp}(c\mid S)=g_c(S),
 \label{eq:history}
\end{align}
where $[x]_+=\max(x,0)$. Thus $g_c$ and $\delta_c$ are different: the
former includes the direct score.

Define the weighted triple count for distinct items
\begin{equation}
 T_{cst}=\sum_b v_b B_{bc}B_{bs}B_{bt}.
\end{equation}
Then
\begin{equation}
 \delta_c(S)=\beta\sum_{\{s,t\}\subseteq S}(a_s+a_t)T_{cst}.
 \label{eq:triple}
\end{equation}
To verify this, a bundle with $k$ query items contributes each seed weight
once for each of the other $k-1$ query items. Grouping the ordered
contributions into unordered pairs gives $a_s+a_t$. For zero overlap,
$o_b=0$, so the same expansion remains valid.

The raw $g_c$ is therefore quadratic in the binary query indicators, with
coefficients tied to observed counts rather than fitted separately. It does
not span every quadratic function. Nor does the whole final pipeline remain
quadratic: query standardization and, in other selected configurations, the
Rao denominator introduce other nonlinear dependence.

\subsection{Category relations and repeated-category penalty}

Let $\ell_i\in\{1,\ldots,L\}$ be the supplied category label of item $i$.
The membership matrix $A\in\{0,1\}^{n\times L}$ has
$A_{i\ell}=\ind\{\ell_i=\ell\}$. Define
\begin{equation}
 H=BA,\qquad h_\ell=\sum_b H_{b\ell},\qquad H_+=\sum_\ell h_\ell,
 \qquad C^{\rm cat}=\off(H^\top H).
\end{equation}
$H_{b\ell}$ counts items of category $\ell$ in bundle $b$;
$\off$ sets the diagonal of a square matrix to zero. For a positive
off-diagonal count, define
\begin{equation}
 P^{\rm cat}_{\ell r}
   =\max\left(0,\log\frac{C^{\rm cat}_{\ell r}H_+}{h_\ell h_r}\right),
\end{equation}
and set it to zero when the count is zero or $\ell=r$. Let
\begin{equation}
 \eta_\ell=
 \frac{\sum_b\ind\{H_{b\ell}\geq2\}+0.5}
      {\sum_b\ind\{H_{b\ell}\geq1\}+1}.
\end{equation}
This is a smoothed empirical repetition rate, not a supplied semantic rule.
The fashion category score is
\begin{equation}
 r_{\rm cat}(c\mid S)=\sum_{s\in S}
 \left[P^{\rm cat}_{\ell_s\ell_c}
 -\frac12\ind\{\ell_s=\ell_c\}\ind\{\eta_{\ell_s}<0.15\}
                      (1-\eta_{\ell_s})\right].
\end{equation}
It pools category-pair observations and penalizes repetition of categories
that rarely repeat in training. All candidates in one category receive the
same value for a fixed query.

\subsection{Descriptions combined with historical-neighbor descriptions}

For each available description or image feature source $u$, let
$x_i^{(u)}$ be its fixed, normalized input vector. These are supplied
representations; this scorer does not train their upstream encoders.
For each category $\ell$, define its arithmetic mean
\begin{equation}
 \bar x_\ell^{(u)}=
 \frac{1}{|\{i:\ell_i=\ell\}|}
       \sum_{i:\ell_i=\ell}x_i^{(u)}.
\end{equation}
Subtract this mean from each item's vector and normalize the result to unit
length when nonzero. Concatenate these residual vectors across available
sources and normalize again, producing row vectors $e_i$. Let $E$ be the
matrix with rows $e_i^\top$.

Define $D=\diag(a_i)$ and the symmetric historical propagation matrix
\begin{equation}
 L_B=\off\left[D B^\top
       \diag\big((d_b+1)^{-1}\big) B D\right].
\end{equation}
This is an adjacency operator, not a Laplacian or a row-stochastic transition
matrix. Let $\rownorm$ normalize each nonzero row to unit Euclidean length.
The bridge representation and score are
\begin{equation}
 W=\rownorm\left(E+\tfrac12\rownorm(L_BE)\right),\qquad
 r_{\rm bridge}(c\mid S)=w_c^\top\frac1{|S|}\sum_{s\in S}w_s,
\end{equation}
where $w_i^\top$ is row $i$ of $W$. It compares an item's own residual
description together with an aggregate of its neighbors' residual
descriptions. This explains why the bridge includes historical information.

\subsection{Co-occurrence profiles compressed by SVD}

Define the item-pair counts and their endpoint totals
\begin{equation}
 C=\off(B^\top B),\qquad e_i^{\rm cnt}=\sum_j C_{ij},\qquad
 C_+=\sum_{ij}C_{ij}.
\end{equation}
The endpoint totals differ from item frequency $f_i$: appearances in larger
bundles contribute more partners. For $C_{ij}>0$, define the positive
pointwise mutual information matrix
\begin{equation}
 P_{ij}=\max\left(0,\log\frac{C_{ij}C_+}
                                  {e_i^{\rm cnt}e_j^{\rm cnt}}\right),
\end{equation}
and set $P_{ij}=0$ otherwise. Let $V_r$ contain retained right singular
vectors in a rank-$r$ SVD of $P$, and form the item representation
$X_{\rm svd}=PV_r=U_r\Sigma_r$ for an exact decomposition. The
implementation uses a truncated numerical approximation. Normalize the
rows of $X_{\rm svd}$ to obtain $u_i$. Then
\begin{equation}
 r_{\rm csvd}(c\mid S)=u_c^\top\frac1{|S|}\sum_{s\in S}u_s.
\end{equation}
This compares compressed association profiles. It does not directly
reconstruct the PPMI entry between each candidate and seed. Orthonormal
singular-vector columns do not imply unit-length rows of $U_r\Sigma_r$,
which is why row normalization is a separate operation. The representation
principle has explicit precedents~\cite{levy2015,afchar2023}.

\subsection{Popularity and selection}

The remaining raw score is $r_{\rm pop}(c\mid S)=\log(1+f_c)$, independent
of the query before standardization. Apply~\eqref{eq:standardize} to each
of the five raw columns, then use~\eqref{eq:dense}.
The displayed constants describe the frozen recipe; these formulas do not
prove that those constants or exactly these five columns are optimal.

\section{A sufficient aggregate under correlated normal measurements}
\label{app:normal}

This appendix gives a conditional theorem connecting several correlated
measurements to a single aggregate and to equation~\eqref{eq:core}. It is
an explanatory model, not a distributional claim about the current raw
scores or about binary hidden-item indicators.

For one candidate and query, let $t$ be an unknown real compatibility
quantity. Condition on the information used to construct a fixed fallback
$F$ and on frequency $f>0$. Suppose a $p$-dimensional vector of
bias-corrected measurements $\mathbf u$ obeys
\begin{equation}
 \mathbf u=\boldsymbol\ell\,t+\boldsymbol\epsilon,\qquad
 \boldsymbol\epsilon\mid t,F,f\sim\Normal(0,\Sigma/f),\qquad
 t\mid F,f\sim\Normal(F,\tau^2).
 \label{eq:normalmodel}
\end{equation}
The response vector $\boldsymbol\ell$, positive-definite covariance
$\Sigma$, and prior variance $\tau^2>0$ are known within this conditional
model. Error correlations are allowed. Define
\begin{equation}
 I=\boldsymbol\ell^\top\Sigma^{-1}\boldsymbol\ell>0,\qquad
 y=\frac{\boldsymbol\ell^\top\Sigma^{-1}\mathbf u}{I},\qquad
 \kappa=\frac1{\tau^2 I}.
\end{equation}

\begin{proposition}[Conditional sufficient aggregate]
For known $\boldsymbol\ell,\Sigma,f$, $y$ is sufficient for $t$ in this
likelihood, and
\begin{equation}
 \E[t\mid\mathbf u,F,f]
    =F+\frac{f}{f+\kappa}(y-F).
\end{equation}
\end{proposition}
\begin{proof}
The part of the negative log-likelihood depending on $t$ is
\begin{equation}
 \frac f2(\mathbf u-\boldsymbol\ell t)^\top
       \Sigma^{-1}(\mathbf u-\boldsymbol\ell t)
   =\frac{fI}{2}(t-y)^2+\text{constant in }t.
\end{equation}
The data enter through $y$. Adding the prior penalty
$(t-F)^2/(2\tau^2)$ and completing the square gives the posterior mean.
This is standard normal conjugacy~\cite{murphy2007}.
\end{proof}

The theorem permits summarizing a specified measurement vector. It does
not say that a new measurement can never help. Suppose a new measurement
has response coefficient $\ell_{\rm new}$, error variance $v/f$, and
cross-covariance $k/f$ with the old errors. Assume the expanded errors are
jointly normal with a positive-definite covariance and remain independent
of $t$ conditional on $F,f$. Define
\begin{equation}
 \ell_\perp=\ell_{\rm new}-k^\top\Sigma^{-1}\boldsymbol\ell,
 \qquad v_\perp=v-k^\top\Sigma^{-1}k>0.
\end{equation}
Subtracting $k^\top\Sigma^{-1}\mathbf u$ from the new measurement
removes the part predictable from old errors. The remainder has response
$\ell_\perp$ to $t$ and independent normal error variance $v_\perp/f$.
Therefore the aggregate precision factor becomes
\begin{equation}
 \boxed{\ I_{\rm new}=I+\frac{\ell_\perp^2}{v_\perp}.\ }
\end{equation}
The additional measurement contributes no information exactly when
$\ell_\perp=0$ under this model. A nearly duplicated measurement can
still help if its remaining part measures $t$ accurately. Exact duplicates
give a singular expanded covariance and should be removed before using
these inverse formulas.

For current BundleFuse, the required assumptions are unverified. Its grid
weights are not these covariance-adjusted coefficients. Its full score and
fallback share inputs. A historical sum is not automatically a sample mean
whose error variance is proportional to $1/f$. The distributional
parameters are not known. A single global $\kappa$ would additionally
require $\tau^2 I$ to be constant across candidate/query conditions.
The negative diagnostic slopes in Section~\ref{sec:claims} reinforce the
need to keep this conditional theorem separate from an unconditional claim
about the implemented scores.

\section{Evidence and reproducibility}
\label{app:evidence}

The mathematical checks are deliberately small and synthetic. They neither
fit new benchmark models nor alter existing checkpoints. Run from the
repository root:
\begin{quote}\ttfamily\small
PYTHONPATH=. python scripts/check\_score\_sufficiency.py
\end{quote}
The retained result is
\path{results/algebraic_diagnostics/score_sufficiency.json}.
It records source hashes and the following checks:
\begin{itemize}
\item The $q=2$ construction: 8 items and 8 bundles per corpus.
\item The $q=21$ construction: 84 items and 9,261 bundles per corpus.
\item Equal frequency, pair and triple counts; equal fashion score vectors
and frozen dense fusion to floating-point tolerance; different four-item
completion counts.
\item The precision-increment identity in 100 positive-definite covariance
cases, with maximum numerical discrepancy $2.78\times10^{-15}$.
\item The expected normalized-relevance ranking rule and additional-utility
bound in 100 Recall and 100 NDCG cases, enumerating all 24 rankings of
four candidates. Maximum discrepancy from the optimal sorted utility is
$2.23\times10^{-16}$.
\end{itemize}
The checks illustrate and audit the derivations; the proofs, rather than
the finite numerical cases, establish the propositions.

The frozen recipes and variance diagnostics are stored in
\path{results/publication_mechanisms/pog/full.json},
\path{results/publication_mechanisms/pog_dense/full.json}, and
\path{results/publication_mechanisms/spotify/full.json}.
The completed controlled studies are recorded under
\path{results/novelty_attribution/} and
\path{results/amplified_fallback/}.
Those experiments support the empirical statements in
Section~\ref{sec:claims}; they are not the independent audit required by
Section~\ref{sec:stopping}.

\clearpage
\begin{thebibliography}{9}
\small
\bibitem{ravikumar2011}
P. Ravikumar, A. Tewari and E. Yang.
\emph{On NDCG Consistency of Listwise Ranking Methods}. AISTATS, 2011.
\url{https://proceedings.mlr.press/v15/ravikumar11a.html}.

\bibitem{banerjee2005}
A. Banerjee, X. Guo and H. Wang.
\emph{On the Optimality of Conditional Expectation as a Bregman Predictor}.
IEEE Transactions on Information Theory 51(7), 2664--2669, 2005.
Author preprint: \url{https://www.dam.brown.edu/people/huiwang/research/bgw.pdf}.

\bibitem{fukumizu2004}
K. Fukumizu, F. R. Bach and M. I. Jordan.
\emph{Dimensionality Reduction for Supervised Learning with Reproducing
Kernel Hilbert Spaces}. JMLR 5, 73--99, 2004.
\url{https://jmlr.org/papers/v5/fukumizu04a.html}.

\bibitem{hoeffding1963}
W. Hoeffding.
\emph{Probability Inequalities for Sums of Bounded Random Variables}.
JASA 58(301), 13--30, 1963.
\url{https://doi.org/10.1080/01621459.1963.10500830}.

\bibitem{murphy2007}
K. P. Murphy.
\emph{Conjugate Bayesian Analysis of the Gaussian Distribution}. 2007.
\url{https://www.cs.ubc.ca/~murphyk/Papers/bayesGauss.pdf}.

\bibitem{levy2015}
O. Levy, Y. Goldberg and I. Dagan.
\emph{Improving Distributional Similarity with Lessons Learned from Word
Embeddings}. TACL 3, 211--225, 2015.
\url{https://aclanthology.org/Q15-1016/}.

\bibitem{afchar2023}
D. Afchar, R. Hennequin and V. Guigue.
\emph{Of Spiky SVDs and Music Recommendation}. RecSys, 2023.
\url{https://arxiv.org/abs/2307.01212}.

\bibitem{coheat2024}
H. Jeon, J. Lee, J. Yun and U. Kang.
\emph{Cold-start Bundle Recommendation via Popularity-based Coalescence
and Curriculum Heating}. The Web Conference, 2024.
\url{https://arxiv.org/abs/2310.03813}.

\bibitem{ramen2025}
\emph{RaMen: Multi-Strategy Multi-Modal Learning for Bundle Construction}.
2025. \url{https://arxiv.org/abs/2507.14361}.
\end{thebibliography}
\end{document}
